\section*{Complete author Section 6.1 proof}
Here, we show that if the truncation of each $(1+k^2x)^{i/k}$, for $i\in[k-1]$, has small constant-free circuit, then one can factor $n$ in $\mathrm{poly}(\log n)$ time, in the non-uniform setting. This would readily imply the first part of Theorem 8.
\paragraph{Proof of Theorem 32.} Let, $(1+k^2x)^{i/k}=\sum_{d\ge0}C_{d,i}x^d\in\mathbb Z[[x]]$, where the coefficient $C_{d,i}$ of $x^d$ is equal to $\pm k^d(-i)(k-i)(2k-i)\cdots((d-1)k-i)/d!$. We see that the product of all $C_{d,i}$ is equal to:
\[\prod_{i\in[k-1]}C_{d,i}=\pm\frac{k^{(k-1)d}(k-1)!(dk)!}{(d!)^kk^d(kd-1)(kd-2)\cdots(kd-(k-1))}.\]
The assumption $\tau(\mathrm{trunc}((1+k^2x)^{i/k},d))=O(\log^cd)$ implies that $\tau(C_{d,i})=O(\log^cd)$ (just by subtracting two consecutive truncations and substituting $x=1$). This further implies that $\tau(\prod_{i\in[k-1]}C_{d,i})=O(\log^cd)$. Let us define, for any $d\ge1$,
\[N(d,k):=\frac{k^{(k-2)d}(dk)!}{(d!)^k}.\]
We first argue that $N(d,k)\in\mathbb N$. This follows from the fact that $N(d,k)=\prod_{i\in[k-1]}C_{d,i}\cdot(kd-1)\cdots(kd-(k-1))/(k-1)!$, and $(k-1)!$ must divide $(kd-1)\cdots(kd-(k-1))$, by Fact 74. Further, since $k$ is constant, it implies that $\tau(N(d,k))=O(\log^cd)$ (because the extra term has trivial $O(\log d)$-complexity).
Now, we describe how to find a non-trivial factor of a given integer $n$. We assume that all the primes dividing $n$ are larger than $k$; otherwise we can remove all the prime factors smaller than $k+1$ (since $k$ is a constant).
The idea is to first find a positive integer $t$ such that all the primes dividing $n$ are in the interval $[t+1,tk]$, by using an iterative algorithm; if such a $t$ does not exist we would have already found a non-trivial factor of $n$ (by the algorithm). As an invariant, we maintain an integer $m$ such that all the prime divisors of $n$ are greater than $m$. We start with $m=k$ and compute $\gcd(N(m,k),n)$ at each iteration. Since all the primes dividing $n$ are greater than $m$ (by assumption), we get that $\gcd(N(m,k),n)=\gcd((mk)!,n)$. If the $\gcd((mk)!,n)\ne1,n$, we must have already found a non-trivial factor of $n$ and we are done. Otherwise, we can have two cases: either (i) $\gcd((mk)!,n)=1$, or (ii) $\gcd((mk)!,n)=n$.
If $\gcd((mk)!,n)=1$ then we set $m\leftarrow mk$ and continue (because in this case all the primes dividing $n$ must be greater than $mk$). Otherwise we have $\gcd((mk)!,n)=n$, and hence, all the primes dividing $n$ are in the interval $[m+1,mk]$ and we stop with $t\leftarrow m$. We know that $t\le\lceil n/k\rceil$ and this uses at most $\log_kn=\log n$ iterations. So, this step has given us an integer $t$ such that all the primes dividing $n$ are in the interval $[t+1,tk]$, and the time taken is $\mathrm{poly}(\log n)$, due to only $\log n$ many iterations and each step takes $\mathrm{poly}(\log n)$-time due to the fact that $\tau(N(d,k))=O(\log^cd)$ implies gcd computation can be done in $\mathrm{poly}(\log n)$ (by euclidean algorithm).
Once, we know that all the primes are in an interval of the form $[t+1,tk]$, we now try to reduce the length of it to $k$ so that, we can simply brute force to get a factor of $n$, otherwise of course our algorithm would already find a factor. The length reduction part is similar to binary search algorithm that we describe below.
To find a positive integer $s$ such that all the primes dividing $n$ are in the interval $[sk+1,(s+1)k]$ (Or we find a non-trivial factor of $n$), again we use an iterative algorithm. As an invariant, we maintain two positive integers $a,b$ such that all the prime divisors of $n$ are in the interval $[ak+1,bk]$. We start with $a=1,b=t$. Our invariant is trivially true at the start.
At each iteration, we set $c=\lceil(a+b)/2\rceil$ and compute $\gcd(N(c,k),n)$. Since $c\le t$ and all the prime divisors of $n$ are larger than $t$, we get that $\gcd(N(c,k),n)=\gcd((ck)!,n)$. Again, we argue in the same way as before. If the gcd is $\ne1,n$, we have already found a non-trivial factor of $n$ and we are done. Otherwise, we have two cases: either (i) $\gcd((ck)!,n)=1$, or (ii) $\gcd((ck)!,n)=n$.
If $\gcd((ck)!,n)=1$ then it is clear that all the primes dividing $n$ are in the interval $[ck+1,bk]$ and hence we set $a\leftarrow c,b\leftarrow b$. If $\gcd((ck)!,n)=n$ then they all the primes dividing $n$ are in the interval $[ak+1,ck]$ and hence we set $a\leftarrow a,b\leftarrow c$. This will terminate when $b-a\le1$. Hence we find the desired positive integer $s$. This uses at most $\log t=\log n$ iterations.
Now we just need to search for the prime divisors of $n$ in the interval $[sk+1,(s+1)k]$ (an interval of constant length). Now, we brute force to finally find a non-trivial factor of $n$. Similarly, this step also takes $\mathrm{poly}(\log n)$ as each gcd computation takes $\mathrm{poly}(\log n)$ time. So, we have successfully found a non-trivial factor of $n$ by the end of this process, repeating this, we can get all the factors in $\mathrm{poly}(\log n)$-time and we are done.\qed
We also refer to Algorithm 1 in Appendix I, reproduced below.
\section*{Complete author Section 6.2 proof}
In this section, we show that easiness of truncation of $(1+k^2x)^{i/k}$ shows that a multiple of $n!$ must be easy. Note that, this may not imply that $n!$ is easy, however, from complexity-theoretic point-of-view, it is believed to be hard because of non-trivial implications. Shub \& Smale [44] proved: If $n!$ is ultimately hard to compute, then $\mathrm P\ne\mathrm{NP}$ over the field of complex numbers. Here, the computation is over Blum-Shub-Smale (BSS) model and can use complex numbers in the algorithm. In fact, a stronger version (known as $\tau$-conjecture) connects $z(f)$, distinct integer roots of $f$ with $\tau(f)$. Recently, [16] showed that a similar conjecture, in the SOS-model, would in fact imply explicit constructions of rigid matrices \& $\mathrm{VP}\ne\mathrm{VNP}$. For similar related works, we refer to [25,27].
\paragraph{Proof of Theorem 35.} Let, $(1+k^2x)^{i/k}=\sum_{d\ge0}C_{d,i}x^d\in\mathbb Z[[x]]$. From the hypothesis, it follows that there exists $c$ such that $\tau(C_{d,i})\le\log^cd$, for each $i\in[k-1]$ (subtract two consecutive terms and substitute $x=1$). Further, from the proof in Section 6.1 (and following the same notation), we know that
\[\prod_{i\in[k-1]}C_{d,i}=\pm\frac{k^{(k-2)d}(k-1)!(dk)!}{(d!)^k(kd-1)(kd-2)\cdots(kd-(k-1))}.\]
Let us define, $a(d,k):=(dk)!/(d!)^k$. Note that, $a(d,k)\in\mathbb N$ (it is the multinomial coefficient $\binom{dk}{d,\ldots,d}$). Further, $k^{(k-2)d}a(d,k)=\prod_{i\in[k-1]}C_{d,i}\cdot(kd-1)\cdots(kd-(k-1))/(k-1)!$, and $(k-1)!$ must divide $(kd-1)\cdots(kd-(k-1))$, by Fact 74. As $k$ is constant, each $kd-i$ can be computed in $O(\log d)$-time trivially. Further, $\tau(\prod_{i\in[k-1]}C_{d,i})\le O(\log^cd)$. As $\tau$ is additive over multiplication, it follows that
\[\tau(k^{(k-2)d}a(d,k))\le O(\log^cd)\implies\tau_1(a(d,k))\le O(\log^cd).\]
Now we recurse by noticing the following trivial identity that $n!=n!/(\lfloor n/k\rfloor)!)^k\cdot((\lfloor n/k\rfloor)!)^k$. We know by the above relation on $a(d,k)$ (and replacing $d:=\lfloor n/k\rfloor$ for some integer $n$) that
\[\tau_1\left(\frac{(k\lfloor n/k\rfloor)!}{(\lfloor n/k\rfloor)!^k}\right)\le O(\log^cn).\]
Further, any integer $n$ can be written as $n=k\lfloor n/k\rfloor+j$ for some $j\le k-1$. Note that $k\lfloor n/k\rfloor+j$ has complexity at most $\log n$ for each $j\in[k-1]$. So, multiplying $k\lfloor n/k\rfloor+j$ for $j\in[k-1]$, it is straightforward to deduce that
\[\tau_1\left(\frac{n!}{(\lfloor n/k\rfloor!)^k}\right)\le O(\log^cn).\tag{2}\]
As, $n!=n!/(\lfloor n/k\rfloor!)^k\cdot((\lfloor n/k\rfloor)!)^k$, and $\tau_1((\lfloor n/k\rfloor!)^k)\le\tau_1(\lfloor n/k\rfloor!)+O(1)$; use Equation (2):
\begin{align*}\tau_1(n!)&\le\tau_1(\lfloor n/k\rfloor!)+O(\log^cn)+O(1)\\&\le\tau_1(\lfloor n/k^2\rfloor!)+O(\log^cn)+O(\log^cn)+O(1)\\&\ \vdots\\&\le\log_kn\cdot O(\log^cn)=O(\log^{c+1}n).\end{align*}
Therefore, $(n!)$ is ultimately easy to compute, as we wanted.\qed
